%% 第四章

\chapter{方法设计与实现}
\chaptermark{方法设计与实现}
\section{基于3DGS的三维重建流程}
在前文分析的基础上可以看出，深空场景三维重建与新视角合成面临远区域细节弱化、高斯形态不稳定以及跨视角外观不一致等问题。针对上述问题，本文在原始3DGS框架基础上，分别引入远区域增强模块FAME、高斯形态稳定模块RVS以及可选的外观一致性模块，以构建适用于深空场景的三维重建与新视角合成方法。为便于说明各部分之间的关系，图4-1给出了本文方法的整体流程。

\begin{figure}[H]
    \centering
    \includegraphics[width=1.00\textwidth]{pipeline.jpg}
    \caption{本文方法整体流程图}\label{fig:4-1}
\end{figure}
如图4-1所示，本文方法整体流程可以分为输入训练视图、初始化与辅助准备、三维高斯重建、改进模块约束以及新视角输出与评价五个阶段。

首先，在输入阶段，系统读取多视角训练图像作为场景观测数据。随后，在初始化与辅助准备阶段，利用COLMAP估计的相机位姿和稀疏点云完成3DGS场景的初始构建；同时，对训练图像进行单目深度估计，并根据深度排序结果提取远区域\textnormal{mask}，为后续FAME模块提供辅助信息。

在3DGS重建阶段，本文仍以高斯的位置、尺度、旋转、不透明度及颜色特征作为优化对象，并通过训练视角渲染结果与真实图像之间的重建误差进行基础优化。

在此基础上，本文进一步引入三类改进模块。其一，FAME模块根据训练视角渲染图像、真实训练图像以及远区域\textnormal{mask}计算远区域加权重建损失，从而增强模型对远处区域的关注。其二，RVS模块通过对高斯形态进行稳定化约束，抑制针状或异常退化的高斯基元，并提升三维高斯空间表示的稳定性，从而减少异常高斯形态对空间结构表达和跨视角投影关系的干扰，进一步约束场景表示的几何一致性。其三，外观一致性模块作为可选模块，在需要处理复杂光照扰动和外观变化时启用。该模块结合高斯位置编码、固有外观特征、参考图像采样特征以及视角方向信息，预测与当前视角一致的高斯颜色，并额外引入采样坐标正则项以约束采样过程稳定。

最终，模型输出新视角合成结果，并采用PSNR、SSIM和LPIPS对重建质量进行评价。图中同时给出了本文方法的总损失构成，其中基础重建损失与FAME损失、RVS损失共同参与优化，而外观一致性项仅在对应模块启用时加入总目标函数。后续各节将围绕图4-1中的关键组成部分展开详细说明。

\section{深度辅助的远区域划分方法}
FAME的核心目的是增强远景区域，但“远景”本身并不是一个固定绝对深度阈值。不同深空场景的尺度、相机轨迹和深度分布差异较大，如果直接采用固定深度值划分远景区域，可能导致某些场景中的远景区域过小，而另一些场景中的远景区域过大，从而影响远景增强约束的稳定性和可比性。

首先读取单目深度估计结果，并在每个视角内对有效深度像素按照深度值进行排序。随后，根据预设远区域比例$r_{\mathrm{far}}$，选取深度分布中位于远端的像素集合构造远区域\textnormal{mask}。该策略避免了对固定绝对深度阈值的依赖，能够根据不同视角的相对深度分布自适应确定远区域范围。

给定辅助深度图$D$和远区域比例$r$，远区域\textnormal{mask}可由深度分位数确定：

\begin{equation}
    M_{\mathrm{far}}(u,v)=\mathbf{1}\left(D(u,v)\ge Q_{1-r}(D_{\mathrm{valid}})\right)
\end{equation}
其中，$Q_{1-r}$表示有效深度集合的1-r分位数，$M_{\mathrm{far}}(u,v)=1$表示该像素被划分为远景区域。本文使用辅助深度结果的目的是根据图像中的相对深度分布确定远区域监督范围。具体而言，当某一像素的深度值大于或等于该视角有效深度集合的$1-r_{\mathrm{far}}$分位数时，说明该像素位于该视角中相对较远的区域，因此被纳入远区域\textnormal{mask}。由此，$r_{\mathrm{far}}$成为控制远景增强覆盖范围的关键超参数：当取值过小时，参与增强的远景像素数量不足；当取值过大时，较多中近景区域也会被纳入远景增强范围，可能削弱方法的针对性。

为了系统确定合理取值，本文从数据集中在固定随机种子下选取10个场景，并设置$r_{\mathrm{far}}$=0.05,0.10,0.15,...,0.50共10个候选值。每个场景按照候选值分别训练10次，共得到100组训练、渲染和评价结果。为避免不同场景之间难度差异对结论造成干扰，本文在每个场景内部对10个$r_{\mathrm{far}}$的结果进行归一化后再统计，从而更准确地反映不同比例本身带来的性能变化。

综合评分同时考虑远景区域渲染质量、全图渲染质量以及辅助结构稳定性指标。其中，远景区域指标用于判断增强目标是否真正改善，全图指标用于约束最终新视角合成质量，辅助结构稳定性指标用于避免只提升图像指标而损害结构一致性。不同远景比例下的归一化综合评分结果如图4-2所示。

\begin{figure}[H]
    \centering
    \includegraphics[width=0.86\textwidth]{不同远景区域比例下的归一化综合得分.png}
    \caption{不同远景区域比例下的归一化综合评分}\label{fig:4-2}
\end{figure}
\vspace{-\baselineskip}
\begin{table}[H]
    \fontsize{10.5pt}{13.125pt}\selectfont
    \caption{不同远景区域比例的关键综合得分}
    \centering\label{tab:4-1}
    \begin{adjustbox}{max width=\textwidth}
    \begin{tabular}{ccccc}
        \toprule
        ratio & far score & global score & geom score & weighted score \\
        \midrule
        0.20 & 0.3208 & 0.0764 & -0.3554 & 0.1583 \\
        0.25 & 0.2943 & 0.2383 & -0.2097 & 0.2047 \\
        0.30 & 0.3392 & 0.3539 & -0.0723 & 0.2811 \\
        0.35 & 0.2259 & 0.5003 & 0.0161 & 0.2631 \\
        0.40 & 0.0587 & 0.5957 & 0.4766 & 0.2556 \\
        0.45 & -0.0072 & 0.8400 & 0.2286 & 0.2400 \\
        0.50 & -0.0242 & 0.8803 & 0.3904 & 0.2641 \\
        \bottomrule
    \end{tabular}
    \end{adjustbox}
\end{table}
由图4-2和表4-1可知，$r_{\mathrm{far}}$=0.30取得最高的weighted score，其值为0.2811，说明该比例在多指标综合评价下表现最好。相比0.20和0.25，0.30在保持较高far score的同时，全图渲染质量进一步提升；相比0.35、0.40、0.45和0.50等更大比例，0.30的远景区域得分更高，说明其远景划分更加集中，增强目标更加明确，没有因为远景范围过宽而削弱远景增强的有效性。

从场景稳定性角度看，不同场景对$r_{\mathrm{far}}$的敏感程度并不完全一致，因此还需要观察该比例是否只依赖个别场景取得优势。图4-3展示了不同远景比例在多场景归一化评价中的排名稳定性。可以看到，中等比例区间整体表现更稳定，其中0.30位于远景增强效果较好的区域，同时没有出现明显的跨场景退化现象，说明该取值具有较好的跨场景适应性和鲁棒性。

\begin{figure}[H]
    \centering
    \begin{subfigure}[b]{0.49\linewidth}
        \includegraphics[width=\linewidth]{不同远景区域比例的跨场景排名稳定性_a.png}
        \caption{Top-3频次统计}
    \end{subfigure}
    \begin{subfigure}[b]{0.49\linewidth}
        \includegraphics[width=\linewidth]{不同远景区域比例的跨场景排名稳定性_b.png}
        \caption{平均排名统计}
    \end{subfigure}
    \caption{不同远景区域比例的跨场景排名稳定性}\label{fig:4-3}
\end{figure}
综上，本文最终选择$r_{\mathrm{far}}$=0.30作为远景区域划分比例。该取值并非人为指定的经验常数，而是在候选比例实验中综合远景质量、全图质量和结构稳定性后确定的默认设置。

\section{FAME远区域增强方法}
在确定远景区域划分比例后，FAME模块利用4.2节得到的远景\textnormal{mask}对训练损失进行补充约束。基础图像重建损失会对整幅图像计算误差，而FAME只在远景\textnormal{mask}内计算额外重建误差，使远处区域纹理、目标边缘和背景层次在优化过程中获得更高关注度，从而缓解整体损失容易被近景或高纹理区域主导的问题。

FAME远区域增强损失定义为远区域内渲染图像与真实图像的平均颜色误差：

\begin{equation}
    \mathcal{L}_{\mathrm{FAME}}=\frac{\left\|M_{\mathrm{far}}\odot(\hat I-I)\right\|_1}{3\sum_{u,v}M_{\mathrm{far}}(u,v)+\epsilon}
\end{equation}
加入 FAME 远区域增强项后的训练损失可写为：

\begin{equation}
    \mathcal{L}=\mathcal{L}_{\mathrm{base}}+\lambda_{\mathrm{fame}}\mathcal{L}_{\mathrm{FAME}}
\end{equation}
在训练过程中，系统根据当前视角读取对应的远景\textnormal{mask}，并将其与渲染图像和真实图像对齐；随后仅在\textnormal{mask}约束范围内计算远景重建误差，并与基础重建损失共同参与反向传播。本文后续实验均采用4.2节确定的$r_{\mathrm{far}}$=0.30作为默认设置。

FAME模块的优点在于实现方式直接，能够嵌入原有3DGS训练流程，不需要改变三维高斯表示形式和可微渲染器结构。该模块与D2GS等关注远处欠拟合问题的研究方向具有相似动机\upcite{ref24}，但本文更侧重于利用深度分位数构造可跨场景适配的远景监督范围，并通过多场景、多比例实验确定最终远景划分参数。

\section{有效秩与体积正则方法}
三维高斯的尺度参数决定其在空间中的形状和覆盖范围，是影响场景几何结构稳定性的重要因素。在3DGS训练过程中，每个高斯都会根据图像重建误差不断调整位置、尺度、旋转、不透明度和颜色等参数。当观测视角充分、纹理信息丰富时，高斯尺度通常能够在多视角约束下逐渐收敛到较合理的空间分布。然而，在视角约束不足、纹理较弱或远区域监督信号较弱的情况下，仅依赖图像重建损失可能无法充分约束高斯的三维形态，部分高斯容易出现异常拉伸、过度膨胀或局部尺度突变等现象，从而导致漂浮伪影、模糊边界和局部几何结构不稳定。

从几何角度看，三维高斯的三个尺度方向共同决定其空间形态。如果某一个尺度方向远大于另外两个方向，高斯会表现出明显的细长形态，容易退化为针状结构；如果三个尺度整体过大，则会造成空间覆盖范围过宽，使局部表面边界变得模糊；如果尺度整体过小，则可能导致场景表示碎片化，影响区域连续性。因此，仅约束单个尺度值并不能充分描述高斯形态是否合理，需要同时考虑尺度分布的均衡性和整体体积变化。基于此，本文引入有效秩正则和体积正则，从高斯形态分布和空间覆盖范围两个角度对三维高斯进行辅助约束。

对于第$i$个高斯，其三个尺度方向记为$s_{i1}$、$s_{i2}$、$s_{i3}$。为了分析三个尺度方向在整体尺度中的相对占比，首先将尺度平方归一化为概率分布：

\begin{equation}
    p_{ik}=\frac{s_{ik}^2}{\sum_{k=1}^{3}s_{ik}^2+\epsilon},\quad k=1,2,3
\end{equation}
其中，$p_{ik}$表示第$i$个高斯在第$k$个尺度方向上的相对能量占比，$\epsilon$为防止分母为零的稳定项。该归一化过程使不同高斯的绝对尺度大小不会直接影响形态分析，而是重点关注三个尺度方向之间的相对分布关系。当某一方向的尺度占据主导时，概率分布会集中在单一方向；当三个尺度方向较为均衡时，概率分布则更加平滑。

在得到尺度方向的归一化分布后，本文利用熵度量其分布均衡程度，并进一步计算有效秩：

\begin{equation}
    H_i=-\sum_{k=1}^{3}p_{ik}\log(p_{ik}+\epsilon),\quad r_i=\exp(H_i)
\end{equation}
其中，$H_i$表示第$i$个高斯尺度分布的信息熵，$r_i$表示对应的有效秩。有效秩可以理解为高斯在空间中具有显著尺度贡献的有效方向数量。当某一个尺度方向明显占优时，$p_{ik}$分布高度集中，熵值较小，有效秩接近1，说明该高斯更接近针状退化形态；当多个方向具有较均衡的尺度贡献时，熵值增大，有效秩也随之增大，说明该高斯具有更稳定的空间支撑形态。因此，有效秩能够从尺度分布角度反映高斯是否存在严重的单方向退化问题。

为了抑制高斯向针状结构退化，本文构造有效秩正则项：

\begin{equation}
    \mathcal{L}_{\mathrm{erank}}=\frac{1}{N}\sum_{i=1}^{N}\exp\left(-(r_i-1)\right)
\end{equation}
其中，$N$表示参与正则计算的高斯数量。当某个高斯的有效秩接近1时，说明其尺度分布高度集中，对应正则项较大；当有效秩增大时，该项逐渐减小。因此，最小化$\mathcal{L}_{\mathrm{erank}}$可以降低针状高斯在训练中的出现概率，使高斯尺度分布更加稳定。需要注意的是，有效秩正则并不是直接强制所有高斯变为完全各向同性结构，而是作为辅助约束抑制极端退化形态，使高斯在满足图像重建需求的同时保持更合理的空间形态。

有效秩正则主要关注三个尺度方向之间的相对分布，但它不能直接限制高斯整体覆盖范围。例如，一个高斯即使三个方向比例较为均衡，也可能整体尺度过大，从而覆盖过宽区域并造成局部细节模糊；相反，如果整体尺度过小，也可能导致场景表示过于碎片化。因此，本文进一步引入体积正则，对高斯的整体空间覆盖范围进行约束。第$i$个高斯的体积可近似表示为三个尺度方向的乘积，体积正则项定义为：

\begin{equation}
    v_i=s_{i1}s_{i2}s_{i3},\quad \mathcal{L}_{\mathrm{vol}}=\frac{1}{N}\sum_{i=1}^{N}\left(\frac{v_i}{\mathrm{median}(\{v_i\})+\epsilon}-1\right)^2
\end{equation}
其中，$\mathrm{median}(\{v_i\})$表示当前高斯体积集合的中位数。相比均值，中位数对异常值更加鲁棒，不容易被少量极大或极小体积的高斯影响。通过将单个高斯体积与整体典型体积进行归一化比较，体积正则能够抑制局部高斯的异常膨胀或过度收缩，使场景中高斯的空间覆盖范围更加稳定。

有效秩正则和体积正则分别从不同角度约束高斯形态。有效秩正则关注尺度方向分布，主要用于抑制单方向异常拉伸和针状退化；体积正则关注整体尺度大小，主要用于限制高斯过度膨胀或过度收缩。二者具有互补关系：前者保证高斯形态不过度退化，后者保证高斯覆盖范围不过度偏离整体分布。将二者结合后，可以在不改变3DGS基本渲染框架的前提下，对高斯基元的几何形态施加更稳定的辅助约束。

在整体训练目标中，有效秩正则和体积正则并不直接替代图像重建损失，而是作为辅助几何约束参与优化。若结合基础图像重建损失和FAME远区域增强损失，整体训练目标可表示为：

\begin{equation}
    \mathcal{L}=\mathcal{L}_{\mathrm{base}}+\lambda_{\mathrm{fame}}\mathcal{L}_{\mathrm{FAME}}+\lambda_{\mathrm{erank}}\mathcal{L}_{\mathrm{erank}}+\lambda_{\mathrm{vol}}\mathcal{L}_{\mathrm{vol}}
\end{equation}
其中，$\mathcal{L}_{\mathrm{base}}$为基础图像重建损失，$\mathcal{L}_{\mathrm{FAME}}$为远区域增强损失，$\lambda_{\mathrm{fame}}$、$\lambda_{\mathrm{erank}}$和$\lambda_{\mathrm{vol}}$分别表示不同损失项的权重系数。通过该联合优化方式，模型一方面能够利用图像重建损失保证渲染结果与真实图像一致，另一方面能够通过FAME模块增强远区域监督，并通过有效秩和体积正则提升三维高斯形态的稳定性。

综上，本文引入的有效秩与体积正则主要用于改善3DGS在几何约束不足时可能出现的高斯形态退化问题。有效秩正则和体积正则从高斯基元形态角度约束三维表示，使模型在视角约束不足和弱纹理区域中仍能保持更加稳定的空间支撑关系，从而提升新视角渲染中的结构一致性。

\section{光照扰动与外观建模设计}
深空场景中的图像外观容易受到太阳高度角、阴影分布、曝光状态、白平衡偏移以及成像设备响应差异等因素影响。对于基于多视角图像监督的3DGS方法而言，训练过程通常假设不同视角之间具有相对稳定的光度一致性，即同一三维区域在不同图像中的颜色观测应当保持相近。然而，在实际成像过程中，同一场景结构可能在不同视角下呈现出亮度、颜色和局部阴影差异。如果直接将这些外观不一致的图像用于训练，模型可能难以区分真实几何结构变化和成像条件变化，从而导致颜色偏移、纹理模糊或新视角渲染不稳定等问题。

为了分析模型在复杂光照条件下的外观建模能力，本文设计了光照扰动数据处理流程。该流程在保持场景几何结构、相机位姿以及训练测试划分不变的前提下，对训练图像施加低光、过曝、白平衡偏移和局部亮度变化等扰动，用于模拟深空场景中可能出现的非均匀光照和成像差异。需要说明的是，该组光照扰动数据主要用于外观一致性相关实验分析，不参与其他评价。通过这种方式，可以在几何条件相对可控的情况下，单独观察外观变化对新视角合成结果的影响，并比较不同外观建模策略的鲁棒性。

在外观建模方面，本文设计了基于参考图像的外观分解分支。该分支借鉴Gaussian in the Wild、WildGaussians和Luminance-GS等工作中关于非受控图像集合外观变化建模的思想，核心目标是在显式三维高斯表示中补充对跨视角外观变化的表达能力\upcite{ref18,ref19,ref20}。与3DGS直接依赖球谐系数表示颜色不同，外观建模分支进一步引入参考图像特征、采样特征、三维位置编码和视角方向信息，使模型能够根据当前观察条件预测更适合该视角的高斯颜色。

具体而言，参考图像首先经过轻量化特征提取网络，得到能够表征当前图像光照和颜色状态的外观特征。该特征并不直接表示场景几何，而是用于描述与成像条件相关的动态外观信息，例如整体亮度、局部阴影、颜色偏移和曝光变化等。随后，系统根据高斯中心在参考图像中的投影位置进行特征采样，并结合若干学习采样位置，将图像空间中的外观信息映射到对应的三维高斯上。这样，每个高斯不仅具有自身的静态颜色和几何属性，还能够获得与参考视角相关的动态外观条件。

外观分支将参考图像特征、学习采样特征、三维点位置编码和视角方向进行融合，其中动态外观特征可表示为：

\begin{equation}
    f_i^{\mathrm{dyn}}=\left[F_p(I_{\mathrm{ref}}(\mu_i)),F_1(q_{i1}),\ldots,F_K(q_{iK})\right]
\end{equation}
式中，参考图像表示用于提供外观状态的输入图像，三维高斯中心位置用于确定该高斯在参考图像中的投影位置，参考图像特征提取函数表示从参考图像中提取并采样得到的外观特征；学习采样坐标表示第$i$个高斯对应的第$k$个采样位置，对应采样函数表示该采样点的特征映射结果。动态外观特征用于描述当前参考条件下可能发生变化的颜色和亮度信息。

这里的动态外观特征与3DGS中的静态颜色特征并不等价。静态颜色特征主要用于描述高斯在场景中的固有外观，例如场景材料、岩石纹理和基础颜色；动态外观特征则用于补充由光照、曝光和相机响应变化引起的视角相关外观差异。通过引入该特征，模型不再完全依赖固定球谐颜色参数去解释所有视角中的颜色变化，而是可以根据参考图像提供的外观信息，对当前视角下的颜色进行自适应预测。

每个高斯的预测颜色由融合网络输出，可写为：

\begin{equation}
    \hat c_i=\mathrm{MLP}\left(f_i^{\mathrm{static}},f_i^{\mathrm{dyn}},\gamma(\mu_i),d_i\right)
\end{equation}
式中，$\hat c_i$表示第i个高斯在给定参考外观和观察方向条件下的颜色估计结果，$f_i^{\mathrm{static}}$表示由球谐系数特征扩展得到的静态外观特征，用于描述高斯的固有外观属性，$f_i^{\mathrm{dyn}}$表示由参考图像引入的动态外观特征，用于描述光照、曝光和颜色响应等成像条件变化，$\gamma(\mu_i)$表示位置编码用于提供高斯在三维空间中的位置信息，$d_i$为视角方向用于补充观察方向对颜色响应的影响，多层感知机$\mathrm{MLP}(\cdot)$用于完成特征融合和颜色预测。该融合过程使颜色预测同时受到场景自身属性、参考图像外观状态、空间位置和观察方向的影响，从而增强模型对复杂光照变化的适应能力。

从建模角度看，静态外观特征和动态外观特征分别对应两类不同信息。前者强调场景中相对稳定的外观属性，例如场景材料、岩石纹理和基础颜色；后者强调由成像条件导致的动态变化，例如局部阴影、曝光增强或颜色偏移。位置编码用于提供高斯在三维空间中的位置信息，使网络能够区分不同区域的外观变化；视角方向则用于补充观察方向对颜色响应的影响。通过多种信息的联合输入，外观分支能够对不同视角下的高斯颜色进行更加灵活的建模。

为了避免学习采样位置偏离有效范围，本文进一步引入采样坐标正则项：

\begin{equation}
    \mathcal{L}_{\mathrm{sc}}=\frac{1}{NK}\sum_{i=1}^{N}\sum_{k=1}^{K}\max\left(\|q_{ik}\|-1,0\right)
\end{equation}
式中，N表示参与训练的高斯数量，K表示每个高斯对应的采样点数量，采样坐标表示第i个高斯的第k个采样位置。该正则项用于约束采样坐标处于合理范围内：当采样坐标位于有效范围内时，对应惩罚为零；当采样坐标超出范围时，正则项会对其产生惩罚，从而避免网络学习到无效或不稳定的采样位置。该约束有助于保证动态外观特征的采样过程稳定，并减少由于无效采样导致的颜色预测偏差。

在整体训练过程中，外观建模分支并不改变3DGS的基本几何表示，而是主要作用于颜色预测阶段。3DGS仍然负责维护高斯的位置、尺度、旋转和不透明度等几何属性，外观分支则根据参考图像和视角信息对高斯颜色进行调节。这样可以在保持显式高斯几何结构的同时，提高模型对不同光照和外观条件的适应能力。

外观一致性模块的整体流程如图4-4所示。

\begin{figure}[H]
    \centering
    \includegraphics[width=0.92\textwidth]{外观一致性模块流程示意图.jpg}
    \caption{外观一致性模块流程示意图}\label{fig:4-4}
\end{figure}
图4-4展示了外观一致性模块的整体流程。输入图像首先经过轻量化特征提取网络得到采样特征和投影特征，随后结合高斯点位置、固有外观特征和视角方向构造动态外观特征。融合网络进一步整合静态外观、动态外观、位置编码和视角信息，并通过颜色解码器输出当前视角下的高斯颜色。该流程使模型能够根据参考图像中的外观状态对高斯颜色进行自适应调整，从而缓解复杂光照条件下的跨视角颜色偏差问题。

将外观建模分支与FAME远区域增强、有效秩正则和体积正则联合后，本文最终采用的训练总损失可写为：

\begin{equation}
    \mathcal{L}=\mathcal{L}_{\mathrm{base}}+\lambda_{\mathrm{fame}}\mathcal{L}_{\mathrm{FAME}}+\lambda_{\mathrm{erank}}\mathcal{L}_{\mathrm{erank}}+\lambda_{\mathrm{vol}}\mathcal{L}_{\mathrm{vol}}+\lambda_{\mathrm{sc}}\mathcal{L}_{\mathrm{sc}}
\end{equation}
式中，基础图像重建损失用于维持渲染结果与真实图像的一致性，FAME远区域增强损失用于提升远景区域监督强度，有效秩正则和体积正则分别用于约束高斯尺度分布和整体空间覆盖范围，采样坐标正则用于稳定外观特征采样过程。通过上述联合优化，模型一方面能够增强远区域细节重建和高斯形态稳定性，另一方面能够补充3DGS对跨视角外观变化的表达能力。

综上，本文通过光照扰动数据处理流程和外观建模分支，从实验数据和模型结构两个层面分析复杂光照条件下的新视角合成问题。光照扰动数据用于构造可控的外观变化场景，便于比较模型在低光、过曝、白平衡偏移和局部亮度变化条件下的表现；外观建模分支则通过参考图像特征和动态外观特征补充3DGS的颜色表达能力。该设计有助于提升模型在跨视角外观不一致条件下的渲染稳定性，并为后续外观一致性实验提供方法基础。

\section{本章小结}
本章介绍了本文方法的设计与实现。首先说明了基于3DGS的基础三维重建流程，随后介绍了深度辅助远区域划分方法和FAME远区域增强模块，并进一步阐述了有效秩正则、体积正则以及光照扰动和外观建模设计。上述方法共同构成了面向深空场景三维重建与新视角合成的改进方案。
